% Part: first-order-logic
% Chapter: model-theory
% Section: partial-iso

\documentclass[../../../include/open-logic-section]{subfiles}

\begin{document}

\olfileid[tr]{mod}{bas}{pis}
\section{Kısmi İzomorfizmalar}

\begin{defn}
  İki !!{structure}s~$\Struct M$ ve $\Struct N$ verilsin. $\Struct M$
  yapısından $\Struct N$ yapısına bir \emph{kısmi izomorfizma},
  girdilerini $\Domain M$ içinden, değerlerini $\Domain N$ içinden alan
  ve kendi tanım kümesi üzerinde \olref[iso]{defn:isomorphism}
  tanımındaki izomorfizma koşullarını sağlayan sonlu bir kısmi
  fonksiyon~$p$'dir:
  \begin{enumerate}
  \item $p$ !!{injective};
  \item her !!{constant}~$c$ için $p(\Assign{c}{M})$ tanımlıysa
    $p(\Assign{c}{M})=\Assign{c}{N}$;
  \item her $n$ yerli !!{predicate}~$P$ için $a_1,\dots,a_n$,
    $p$'nin tanım kümesindeyse
    $\langle a_1,\dots,a_n\rangle\in\Assign P M$ ancak ve ancak
    $\langle p(a_1),\dots,p(a_n)\rangle\in\Assign P N$;
  \item her $n$ yerli !!{function}~$f$ için $a_1,\dots,a_n$,
    $p$'nin tanım kümesindeyse
    $p(\Assign f M(a_1,\dots,a_n))=
    \Assign f N(p(a_1),\dots,p(a_n))$.
  \end{enumerate}
  $p$'nin sonlu olması, $\dom p$ kümesinin sonlu olması demektir.
\end{defn}

Boş fonksiyon~$\emptyset$, herhangi iki !!{structure}s arasında her
zaman bir kısmi izomorfizmadır.

\begin{defn}\ollabel{defn:partialisom}
  İki !!{structure}s~$\Struct M$ ve $\Struct N$ arasında aşağıdaki
  \emph{ileri-geri} özelliğini sağlayan, boş olmayan bir kısmi
  izomorfizmalar kümesi~$I$ varsa ve ancak varsa bunların \emph{kısmen
  izomorf} olduğunu söyler ve $\Struct M\iso[p]\Struct N$ yazarız:
  \begin{enumerate}
  \item (\emph{İleri}) Her $p\in I$ ve $a\in\Domain M$ için
    $p\subseteq q$ olan ve $a$'yı tanım kümesinde içeren bir $q\in I$
    vardır;
  \item (\emph{Geri}) Her $p\in I$ ve $b\in\Domain N$ için
    $p\subseteq q$ olan ve $b$'yi görüntü kümesinde içeren bir $q\in I$
    vardır.
  \end{enumerate}
\end{defn}

\begin{thm}\ollabel{thm:p-isom1}
  $\Struct M\iso[p]\Struct N$ ve $\Struct M$ ile $\Struct N$
  !!{enumerable} ise $\Struct M\iso\Struct N$.
\end{thm}

\begin{proof}
  $\Struct M$ ve $\Struct N$ !!{enumerable} olduğundan
  $\Domain M=\{a_0,a_1,\ldots\}$ ve
  $\Domain N=\{b_0,b_1,\ldots\}$ yazalım. Gelişigüzel bir $p_0\in I$
  ile başlayıp artan bir
  $p_0\subseteq p_1\subseteq p_2\subseteq\cdots$ kısmi izomorfizmalar
  dizisini şöyle tanımlarız:
  \begin{enumerate}
  \item $n+1$ tekse, $n=2r$ yazıp ileri özelliğiyle
    $p_n\subseteq p_{n+1}$ olan ve $a_r$'yi tanım kümesinde içeren bir
    $p_{n+1}\in I$ bulun;
  \item $n+1$ çiftse, $n+1=2r$ yazıp geri özelliğiyle
    $p_n\subseteq p_{n+1}$ olan ve $b_r$'yi görüntü kümesinde içeren bir
    $p_{n+1}\in I$ bulun.
  \end{enumerate}
Şimdi
\[
p=\bigcup_{n\geq0}p_n
\]
alırsak $p$, $\Struct M$ ile $\Struct N$ arasında bir izomorfizmadır.
\end{proof}

\begin{prob}
  \olref[mod][bas][pis]{thm:p-isom1} içinde tanımlanan $p$'nin gerçekten
  bir izomorfizma olduğunu ayrıntılı olarak gösterin.
\end{prob}

\begin{thm}\ollabel{thm:p-isom2}
  $\Struct M$ ve $\Struct N$, salt bağıntısal !!a{language} için
  !!{structure}s olsun (yani !!a{language} yalnız !!{predicate}s
  içerip hiçbir !!{function}s veya sabit içermesin). O zaman
  $\Struct M\iso[p]\Struct N$ ise
  $\Struct M\elemequiv\Struct N$.
\end{thm}

\begin{proof}
  !!{formula}s üzerinde tümevarımla şu gösterilir: $a_1,\dots,a_n$ ile
  $b_1,\dots,b_n$ arasında her $a_i$'yi $b_i$'ye eşleyen bir kısmi
  izomorfizma~$p$ varsa ve $s_1(x_i)=a_i$, $s_2(x_i)=b_i$
  ($i=1,\dots,n$) ise $\Sat{M}{!A}[s_1]$ ancak ve ancak
  $\Sat{N}{!A}[s_2]$. $n=0$ durumu
  $\Struct M\elemequiv\Struct N$ sonucunu verir.
\end{proof}

\begin{rem}
!!{function}s bulunuyorsa önceki sonuç yine doğrudur; fakat $p$'nin,
$\Struct M$ yapısının $a_1,\dots,a_n$ tarafından üretilen alt
!!{structure} ifadesi ile $\Struct N$ yapısının $b_1,\dots,b_n$
tarafından üretilen alt !!{structure} ifadesi arasında oluşturduğu
izomorfizmayı ele almak gerekir.
\end{rem}

Önceki sonuç, !!a{formula} içindeki iç içe niceleyicilerin sayısı ile
ilgili kısmi izomorfizmaların kaç kez genişletilebildiği arasında bağ
kurularak aşamalara ``ayrılabilir.''

\begin{defn}
  Her !!{formula}~$!A$ için $!A$'nın \emph{niceleyici derecesi},
  $\QuantRank{!A}\in\Nat$, $!A$ içindeki iç içe niceleyicilerin en
  yüksek sayısı olarak özyinelemeyle tanımlanır. İki
  !!{structure}s~$\Struct M$ ve $\Struct N$, niceleyici derecesi en çok
  $n$ olan bütün tümceler üzerinde uyuşuyorsa \emph{$n$-eşdeğerdir};
  bunu $\Struct M\elemequiv[n]\Struct N$ ile gösteririz.
\end{defn}

\begin{prop}\ollabel{prop:qr-finite}
  $\Lang L$, sonlu sayıda !!{predicate}s ve !!{constant}s içeren,
  hiçbir !!{function}s içermeyen sonlu salt bağıntısal !!a{language}
  olsun. O zaman her $n\in\Nat$ için $\Lang L$ !!{language} içindeki
  niceleyici derecesi en çok $n$ olan birinci dereceden
  !!{sentence}s, mantıksal eşdeğerlik altında yalnızca sonlu sayıdadır.
\end{prop}

\begin{proof}
  $n$ üzerinde tümevarımla.
\end{proof}

\begin{defn}
  !!a{structure}~$\Struct M$ verildiğinde $\Domain M^{<\omega}$,
  $\Domain M$ üzerinde bütün sonlu dizilerin kümesi olsun. Sonlu eleman
  dizileri için $\mathbf a,\mathbf b,\mathbf c,\ldots$ kullanacağız.
  $\mathbf a\in\Domain M^{<\omega}$ ve $a\in\Domain M$ ise
  $\mathbf aa$, $\mathbf a$ ile $a$'nın \emph{art arda eklenmesini}
  gösterir.
\end{defn}

\begin{defn}
  !!{structure}s~$\Struct M$ ve $\Struct N$ verildiğinde, eşit
  uzunluktaki diziler arasında
  $I_n\subseteq\Domain M^{<\omega}\times\Domain N^{<\omega}$
  bağıntılarını $n$ üzerinde özyinelemeyle şöyle tanımlarız:
   \begin{enumerate}
   \item $I_0(\mathbf a,\mathbf b)$ ancak ve ancak $\mathbf a$ ile
     $\mathbf b$, $\Struct M$ ve $\Struct N$ içinde aynı atomik
     !!{formula}s ifadelerini sağlıyorsa; yani $s_1(x_i)=a_i$,
     $s_2(x_i)=b_i$ ve bütün !!{variable}s $x_1,\dots,x_n$ arasında
     bulunan $!A$ atomikse $\Sat{M}{!A}[s_1]$ ancak ve ancak
     $\Sat{N}{!A}[s_2]$ ise.
   \item $I_{n+1}(\mathbf a,\mathbf b)$ ancak ve ancak her
     $a\in\Domain M$ için $I_n(\mathbf aa,\mathbf bb)$ olacak bir
     $b\in\Domain N$ varsa ve bunun tersi de geçerliyse.
   \end{enumerate}
\end{defn}

\begin{defn}
  $\Struct M$ ile $\Struct N$ için
  $I_n(\emptyseq,\emptyseq)$ geçerliyse
  $\Struct M\approx_n\Struct N$ yazarız (burada $\emptyseq$ boş
  dizidir).
\end{defn}

\begin{thm}\ollabel{thm:b-n-f}
  $\Lang L$ salt bağıntısal !!a{language} olsun.
  $I_n(\mathbf a,\mathbf b)$ ise $\QuantRank{!A}\leq n$ olan her
  $!A$ için $\Sat{M}{!A}[\mathbf a]$ ancak ve ancak
  $\Sat{N}{!A}[\mathbf b]$ olur (burada yine $s(x_i)=a_i$ olan herhangi
  bir $s$, $!A$'yı sağlıyorsa $\mathbf a$ da $!A$'yı sağlar). Üstelik
  $\Lang L$ sonluysa bunun tersi de geçerlidir.
\end{thm}

\begin{proof}
  $I_n(\mathbf a,\mathbf b)$ bağıntısının $\mathbf a$ ile $\mathbf b$
  dizilerinin niceleyici derecesi en çok $n$ olan aynı !!{formula}s
  ifadelerini sağladığını gerektirdiği, $!A$ üzerinde kolay bir
  tümevarımla gösterilir. Ters yön için $n$ üzerinde tümevarım yapıp,
  her $n$ için bu niceleyici derecesinde eşdeğer olmayan en fazla sonlu
  sayıda !!{formula}s bulunduğunu güvence altına alan
  \olref{prop:qr-finite} önermesini kullanırız.

  $n=0$ için $\mathbf a$ ile $\mathbf b$'nin aynı niceleyicisiz
  !!{formula}s ifadelerini sağladığı varsayımı, aynı atomik formülleri
  sağladıklarını; dolayısıyla $I_0(\mathbf a,\mathbf b)$ olduğunu verir.

  $n+1$ durumu için $\mathbf a$ ile $\mathbf b$'nin niceleyici derecesi
  en çok $n+1$ olan aynı !!{formula}s ifadelerini sağladığını
  varsayalım. $I_{n+1}(\mathbf a,\mathbf b)$ göstermek için her
  $a\in\Domain M$ için $I_n(\mathbf aa,\mathbf bb)$ olacak bir
  $b\in\Domain N$ bulunduğunu göstermek yeterlidir. Tümevarım
  varsayımı gereği, her $a\in\Domain M$ için $\mathbf aa$ ile
  $\mathbf bb$'nin niceleyici derecesi en çok $n$ olan aynı
  !!{formula}s ifadelerini sağlayacağı bir $b\in\Domain N$ bulunduğunu
  göstermek yine yeterlidir.

  $a\in\Domain M$ verildiğinde $!T^a_n$, $\Struct M$ içinde
  $\mathbf aa$ tarafından sağlanan, derecesi en çok $n$ olan
  $!B(x,\mathbf y)$ !!{formula}s kümesi olsun. Eşdeğerlik altında
  $!T^a_n$ sonlu olduğundan, temsilcilerinin sonlu birleşimini tek bir
  birinci dereceden !!{formula} olarak alabiliriz. O zaman $\mathbf a$,
  niceleyici derecesi en çok $n+1$ olan
  $\lexists[x][!T^a_n(x,\mathbf y)]$ formülünü sağlar. Varsayım gereği
  $\mathbf b$ de $\Struct N$ içinde aynı !!{formula} ifadesini sağlar;
  dolayısıyla $\mathbf bb$'nin $!T^a_n$ formülünü sağladığı bir
  $b\in\Domain N$ vardır. Özellikle $\mathbf bb$, $\mathbf aa$ ile
  niceleyici derecesi en çok $n$ olan aynı !!{formula}s ifadelerini
  sağlar. Benzer biçimde, her $b\in\Domain N$ için $\mathbf aa$ ile
  $\mathbf bb$'nin niceleyici derecesi en çok $n$ olan aynı
  !!{formula}s ifadelerini sağladığı bir $a\in\Domain M$ bulunduğu
  gösterilir; ispat tamamlanır.
\end{proof}

\begin{cor}\ollabel{cor:b-n-f}
  $\Struct M$ ve $\Struct N$, sonlu bir !!{language} içindeki salt
  bağıntısal !!{structure}s ise $\Struct M\approx_n\Struct N$ ancak ve
  ancak $\Struct M\elemequiv[n]\Struct N$. Özellikle
  $\Struct M\elemequiv\Struct N$ ancak ve ancak her $n$ için
  $\Struct M\approx_n\Struct N$.
\end{cor}

\end{document}
